\section{Билет 5. Кинетическая и потенциальная энергии гармонических колебаний. Полная энергия гармонического колебания. Графики.}

\begin{definition}
Кинетическая энергия тела массы $m$, совершающего гармонические колебания, определяется мгновенной скоростью $v(t) = \dot{x}(t) = -A\omega_0 \sin(\omega_0 t + \varphi)$:
\begin{equation}
    T(t) = \frac{m v^2}{2} = \frac{m A^2 \omega_0^2}{2} \sin^2(\omega_0 t + \varphi). \label{eq:kinetic_05}
\end{equation}
Используя тригонометрическое тождество $\sin^2\alpha = \frac{1-\cos 2\alpha}{2}$, запишем:
\begin{equation}
    T(t) = \frac{m A^2 \omega_0^2}{4} \left[1 - \cos(2(\omega_0 t + \varphi))\right].
\end{equation}
Максимальное значение кинетической энергии достигается при прохождении положения равновесия ($x=0$):
\begin{equation}
    T_{\max} = \frac{m A^2 \omega_0^2}{2}.
\end{equation}
\end{definition}

\begin{definition}
Потенциальная энергия упругой (квазиупругой) силы $F = -kx$, отсчитываемая от положения равновесия ($U(0)=0$), равна:
\begin{equation}
    U(x) = \int_0^x kx' dx' = \frac{k x^2}{2} = \frac{k A^2}{2} \cos^2(\omega_0 t + \varphi). \label{eq:potential_05}
\end{equation}
Учитывая, что $k = m \omega_0^2$, получаем:
\begin{equation}
    U(t) = \frac{m A^2 \omega_0^2}{2} \cos^2(\omega_0 t + \varphi) = \frac{m A^2 \omega_0^2}{4} \left[1 + \cos(2(\omega_0 t + \varphi))\right].
\end{equation}
Максимальное значение потенциальной энергии достигается в точках максимального отклонения ($x = \pm A$):
\begin{equation}
    U_{\max} = \frac{k A^2}{2} = \frac{m A^2 \omega_0^2}{2}.
\end{equation}
\end{definition}

\begin{theorem}[Полная энергия гармонического осциллятора]
Полная механическая энергия консервативной системы равна сумме кинетической и потенциальной энергий:
\begin{equation}
    E = T(t) + U(t) = \frac{m A^2 \omega_0^2}{2} \left[\sin^2(\omega_0 t + \varphi) + \cos^2(\omega_0 t + \varphi)\right].
\end{equation}
Так как $\sin^2\alpha + \cos^2\alpha = 1$, полная энергия остаётся постоянной во времени:
\begin{equation}
    E = \frac{m A^2 \omega_0^2}{2} = \frac{k A^2}{2} = \text{const}. \label{eq:total_energy_05}
\end{equation}
\end{theorem}

\begin{property}
Из формул \eqref{eq:kinetic_05} и \eqref{eq:potential_05} следуют ключевые энергетические свойства:
\begin{enumerate}
    \item Кинетическая и потенциальная энергии изменяются с циклической частотой $2\omega_0$, то есть в \textbf{два раза быстрее}, чем смещение и скорость.
    \item Энергии колеблются в \textbf{противофазе}: когда $T$ максимальна, $U=0$, и наоборот. Происходит периодический обмен энергией между кинетической и потенциальной формами без потерь.
    \item Средние за период значения энергий равны: $\langle T \rangle = \langle U \rangle = \frac{E}{2}$.
    \item Полная энергия пропорциональна квадрату амплитуды $A^2$ и квадрату циклической частоты $\omega_0^2$.
\end{enumerate}
\end{property}

\begin{remark}
\textbf{Графики энергий.} На графиках $T(t)$ и $U(t)$ представляют собой «поднятые» синусоиды (квадраты синуса и косинуса), лежащие строго выше оси времени. Их сумма $E(t)$ --- горизонтальная прямая, что наглядно демонстрирует закон сохранения механической энергии. При построении графиков от координаты $x$ зависимость $U(x)$ --- парабола, а $T(x) = E - U(x)$ --- «перевёрнутая» парабола, симметричная относительно оси энергии.
\end{remark}
